%% Section 5: evaluation, comparing the denotational proof with the three
%% operational ones (RQ2, RQ4). Budget: 3.0 pages
%% (notes/paper-outline.md, section 5). Numbers come from loc.md
%% (loc.py --mu) and will change; update them before submission.
%% TODO: setup, accounting rules, Table 1, analysis by kind of pass, full
%% cost, threats to validity.

\section{Evaluation}
\label{sec:evaluation}

\paragraph{Delay is the pass whose denotational proof costs the most.}
The denotational proof of delay takes 524 lines of Agda, against 35 to
139 lines for the same pass in the three operational styles. Most of the difference has three
sources. First, the denotational proof needs two directions, while a
simulation proof gets the backward direction for free from forward
simulation and the determinism of the target. The graph model has an
analogue of determinism, the consistency of denotations, and it does
part of the job: if the source produces a constant, the forward
direction puts that constant in the target, and consistency rules out
any other. But the graph model represents divergence as the empty set,
the absence of evidence, and a forward relation cannot carry absence.
Showing that a diverging program compiles to a diverging program is
therefore a separate argument, and that argument is reflection. (We use
different relations for the two directions, the structural one for
preservation and the approximant one for reflection, because a single
relation parameterized by the direction made the proof longer: on an
earlier version of the compiler, switching from one relation to two cut
the proof of delay from 1,358 to 832 lines.)
Second, reflection needs the finite lists, the batches of calls, and
the lemma that extracts them from a target call. Third, $\beta$-rules in
the graph model hold only under side conditions of nonemptiness and
continuity, which a machine gets for free by taking a step.
